\chapter{Holomorphic BF theory and its boundary current}
\label{chap:zero-level-bf}

A current algebra can lose its central term while its bulk field theory
retains a nondegenerate Batalin--Vilkovisky pairing.  The mechanism is a
change in the differential, with two field components kept separate.
On a complex curve times a real half-line, the quadratic action density is
\[
 \mathcal L_\kappa(a,b)
 =b\wedge d'a+\frac\kappa2a\wedge\partial a,
 \qquad d'=\bar\partial+d_t.
\]
It is defined for smooth fields.  On compactly supported fields the
action is $S_\kappa(a,b)=\int_M^{\mathrm{res}}\mathcal L_\kappa(a,b)$.
Here $a$ is a partial connection in the antiholomorphic and real
directions, and $b$ is a holomorphic one-form.  The pairing between
these fields is independent of $\kappa$.  At $\kappa=0$ the action is
holomorphic BF theory: the field $b$ multiplies the partial curvature
$d'a$.  A compactly supported collar map identifies its quantum
boundary observables with the boundary value of the bulk theory.
A Cauchy contraction then extracts every current field and operator
product from those supported observables.  With the positive unary
differential fixed below, the boundary current has level $-\kappa$.
At zero its singular product vanishes, while its polynomial fields
and all holomorphic jets remain.

\section{The fields, including their odd components}

Let $\Sigma=\C$ with coordinate $z=x+iy$, and let
$M=\Sigma\times\R_{\ge0}$ with coordinate $t$ on the second factor.
Choose the orientation whose induced boundary orientation on
$\Sigma\times\{0\}$ is the complex orientation.
Write $\mathcal A$ for smooth forms generated by $d\bar z,dt$.
All coefficients are complex-valued, and
\[
 d'=d\bar z\wedge\partial_{\bar z}+dt\wedge\partial_t,
 \qquad \partial=dz\wedge\partial_z.
\]
These operators satisfy $(d')^2=\partial^2=d'\partial+\partial d'=0$.
The parameter $\kappa\in\C$ is a scalar of cohomological degree zero.
Statements polynomial in $\kappa$ also hold over $\C[\kappa]$.

The coefficient space of $\alpha$ is $\C e$, and that of $\beta$ is
its dual $\C e^\vee$.  These chosen bases suppress the evaluation
pairing in the scalar formulas.  The fields have two summands,
\[
 E_\kappa=\mathcal A[1]\oplus(\mathcal A\wedge dz)[1].
\]
This is the unrestricted smooth field complex.  Write $E_{\kappa,c}$
for its subcomplex of compactly supported fields; compact support is
required in both the curve and the real direction and may meet the
boundary of $M$.
The notation records the following precise degrees: an ordinary
$p$-form in either summand has cohomological degree $p-1$.
The differential is the actual exterior operator
\[
 Q_\kappa(\alpha,\beta)
   =(d'\alpha,\ d'\beta+\kappa\partial\alpha).
\]
Thus the holomorphic one-form $\beta$ has degree zero before adding
any $d\bar z$ or $dt$.  All tensor substitutions use the Koszul sign
of the displayed cohomological degrees.

For a component of the field complex, the corresponding coordinate
has the opposite degree.  This coordinate degree is the ghost number.
The complete list is
\[
\begin{array}{c|c|c|l}
 \text{component}&\text{form degree}&\text{ghost number}&\text{role}\\\hline
 \alpha^{(0)}&0&1&\text{gauge ghost }c\\
 \alpha^{(1)}=a&1&0&\text{partial connection}\\
 \beta^{(1)}=b&1&0&\text{holomorphic one-form}\\
 \alpha^{(2)}&2&-1&\text{antifield of }b\\
 \beta^{(2)}&2&-1&\text{antifields of }a\\
 \beta^{(3)}&3&-2&\text{antifield of }c .
\end{array}
\]
There is no $\alpha^{(3)}$, because $\mathcal A$ has only two
exterior generators.  Both summands are essential to the BF presentation.

To fix the current normalization, define
\[
 \int^{\mathrm{res}}=\frac1{2\pi i}\int
\]
on compactly supported complex-valued top forms, both on $M$ and on $\Sigma$.
The second integral is ordinary oriented integration.  In contour
notation $\operatorname{Res}_{z=0}z^{-1}dz=1$.
For a compactly supported two-form $\theta$ on $\Sigma$ and a
compactly supported function $\rho$ on the half-line, fiber integration
has the sign
\[
 \int_M^{\mathrm{res}}\theta\wedge\rho(t)dt
 =-\left(\int_0^\infty\rho(t)dt\right)
                  \int_\Sigma^{\mathrm{res}}\theta.
\]
Indeed, the orientation of $M$ is the negative of the product
orientation on $\Sigma\times\R_{\ge0}$.
Changing the common scalar normalization rescales the action and
inverse pairing together.

For homogeneous smooth fields $u=(\alpha,\beta)$ of form degree $p$
and $v=(\alpha',\beta')$, define the local pairing density
\[
 \omega_{\mathrm{loc}}(u,v)=(-1)^{p+1}
       (\alpha\wedge\beta'+\beta\wedge\alpha')_{[3]},
\]
where $[3]$ takes the top-form component.  The orientation identifies
this top form with a complex-valued density.  Its integral
\[
 \omega(u,v)=\int_M^{\mathrm{res}}\omega_{\mathrm{loc}}(u,v)
\]
is defined whenever at least one argument has compact support.
Thus the integrated pairing is defined on
$E_{\kappa,c}\otimes E_\kappa$ and on
$E_\kappa\otimes E_{\kappa,c}$; no integral is assigned to an arbitrary
pair of unrestricted smooth fields.
The local pairing has degree $-1$, is independent of $\kappa$, and is
fiberwise nondegenerate: the two complementary summands pair by
evaluation and wedge product.  Its sign gives graded skew-symmetry.
Indeed, a nonzero pairing has $p+q=3$, so interchanging the forms leaves
their wedge sign unchanged and changes $(-1)^{p+1}$ to its negative.

The distinction is already necessary for fields away from the
boundary.  If $\chi$ is a nonzero real smooth function supported in
$(1,2)$, then $u=(\chi(t)dt,0)$ and
$v=(0,\chi(t)dz\wedge d\bar z)$ have a well-defined local pairing
density, but its integral contains the infinite area of $\C$.
Their support is not compact in $M$.

\section{The BV identity and the chiral boundary condition}

The square of the differential is zero for every $\kappa$:
\[
 Q_\kappa^2(\alpha,\beta)
 =\bigl(0,\kappa(d'\partial+\partial d')\alpha\bigr)=0.
\]
The linear complex is elliptic, including at zero level.  For a
nonzero real covector $\xi$, the symbol of $d'$ is wedge multiplication
by $\xi^{0,1}+\xi_tdt$.  This one-form is nonzero: a real covector
with zero $(0,1)$ part and zero $t$ part is zero.  Contraction with a
dual vector contracts this symbol complex.  Choose that vector with
zero $z$ component.  Its contraction anticommutes with the symbol of
$\partial$, so it also contracts the symbol of $Q_\kappa$.

\begin{proposition}[A nondegenerate BV family]
\label{prop:bf-bv-family}
For every $\kappa\in\C$, the fields, differential and local pairing above
define a free BV theory in the interior of $M$.  The boundary condition
\[
 L=\Omega^{1,\bullet}(\Sigma),\qquad
 \iota^*\alpha=0,\quad \iota^*\beta\in L,
 \qquad \iota:\Sigma\hookrightarrow M,
\]
is a local Lagrangian boundary condition.  The restricted differential
is compatible with the integrated pairing whenever at least one
argument has compact support.
\end{proposition}
\begin{proof}
Ellipticity and nondegeneracy were proved above.  For homogeneous $u$
of form degree $p$, the exterior Leibniz rule gives the local identity
\[
 \begin{split}
 &\omega_{\mathrm{loc}}(Q_\kappa u,v)
              +(-1)^{p-1}\omega_{\mathrm{loc}}(u,Q_\kappa v)\\
 &\quad=(-1)^p\left[d'(\alpha\wedge\beta'+\beta\wedge\alpha')
                  +\kappa\partial(\alpha\wedge\alpha')\right]_{[3]}.
 \end{split}
\]
Assume now that at least one of $u,v$ has compact support.  The
primitives on the right then have compact support, so their Stokes
integrals have no contribution from infinity.  The second term
integrates to zero: its possible primitive is a two-form generated
by $d\bar z,dt$, whose boundary pullback vanishes.  The first term is
an ordinary total derivative, since applying $\partial$ to its
primitive gives holomorphic degree two.  Consequently
\[
 \begin{split}
 &\omega(Q_\kappa u,v)+(-1)^{p-1}\omega(u,Q_\kappa v)\\
 &\quad=(-1)^p\int_\Sigma^{\mathrm{res}}
    \iota^*(\alpha\wedge\beta'+\beta\wedge\alpha').
 \end{split}
\]
Within this support domain the expression vanishes if one argument
is supported away from the boundary, or if both satisfy $L$.

On the boundary, the two Dolbeault summands have equal rank and pair
nondegenerately with one another.  The second summand is isotropic,
so it is Lagrangian.  It is preserved by $Q_\kappa$ because
$\partial\beta=0$.  Its differential is the Dolbeault operator on
holomorphic one-forms, an elliptic complex on $\Sigma$.
The same decomposition tensored with normal forms makes the theory
topological in the normal direction.  When shifts precede this tensor
product, the identification of a normal one-form carries the usual
suspension sign; equivalently one may use the exterior-form formula
for $Q_\kappa$ displayed above.  Thus the boundary condition is a Lagrangian elliptic subcomplex,
and the bulk complex is topological in its normal direction.
\end{proof}

Fix the Hamiltonian convention $\iota_{Q_\kappa}\omega=-\delta S$.
The full graded smooth field complex has the local quadratic density
\[
 S_{\mathrm{loc}}(\Phi)
       =\frac12\omega_{\mathrm{loc}}(\Phi,Q_\kappa\Phi).
\]
For $\Phi\in E_{\kappa,c}$ its integral is the well-defined functional
$S(\Phi)=\int_M^{\mathrm{res}}S_{\mathrm{loc}}(\Phi)
=\frac12\omega(\Phi,Q_\kappa\Phi)$.
Write $E_{\kappa,L,c}=\{\Phi\in E_{\kappa,c}:\iota^*\alpha=0\}$.
With the chiral boundary condition imposed, this functional and its
Hamiltonian identity are taken on that space.
For unrestricted fields, the action is interpreted as a local
density, whose first variation is tested against compactly supported
variations.

On ghost-number-zero components the density is
\[
 \frac12(a\wedge d'b+b\wedge d'a+\kappa a\wedge\partial a)
 =\mathcal L_\kappa(a,b)-\frac12d'(a\wedge b).
\]
For compactly supported fields satisfying $\iota^*a=0$, Stokes'
theorem kills the total derivative, so its integral is the action
stated at the start of the chapter.  The local Euler--Lagrange
equations, valid on the unrestricted smooth field complex, are
\[
 d'a=0,\qquad d'b+\kappa\partial a=0.
\]
Their local gauge symmetry, with smooth gauge parameter $c$ vanishing
on the boundary, is
\[
 \delta_c a=d'c,\qquad \delta_c b=\kappa\partial c.
\]
A compactly supported $c$ also preserves the compactly supported
field space.  The second equation is invariant because
$d'\partial=-\partial d'$.

For any smooth configuration satisfying $L$ and any compactly
supported variation $(\delta a,\delta b)$ satisfying the same boundary
condition, the first variation of the local action integrates to
\[
 \begin{split}
 \int_M^{\mathrm{res}}\delta\mathcal L_\kappa
 &=\int_M^{\mathrm{res}}
       \bigl(\delta b\wedge d'a
              +\delta a\wedge(d'b+\kappa\partial a)\bigr)\\
 &\qquad-\int_\Sigma^{\mathrm{res}}b\wedge\delta a.
 \end{split}
\]
All integrands have compact support because each contains a variation
or its derivative.  Integration by parts is therefore legitimate,
and the boundary term vanishes.  This proves the variational
Hamiltonian identity with its support domain.  The configuration
$a=t\,d\bar z$, $b=dz$ illustrates the remaining distinction: it
satisfies the boundary condition but has a nonintegrable constant
action density.  Its local variational equations are defined even
though it has no integrated action value.

The terms involving ghosts and antifields are fixed by
$S_{\mathrm{loc}}$ on the unrestricted graded fields.  On
$E_{\kappa,L,c}$ the integrated functional satisfies the classical
master equation: the Hamiltonian vector field of $\{S,S\}$ is
$2Q_\kappa^2=0$, and this quadratic functional has no constant term.
Equivalently, the same statement is an identity of local variational
functionals when evaluated against compactly supported variations.

\section{The nonzero-level change of fields and its special fiber}

For $\kappa\ne0$, put
\[
 F_\kappa(\alpha,\beta)=\alpha+\kappa^{-1}\beta.
\]
The target is the full de Rham field complex $\Omega^\bullet(M)[1]$
with differential $d=d'+\partial$ and local pairing density
$\omega_{\mathrm{CS},\kappa,\mathrm{loc}}(u,v)
=(-1)^{p+1}\kappa(u\wedge v)_{[3]}$.
Its integrated pairing is
$\omega_{\mathrm{CS},\kappa}=\int_M^{\mathrm{res}}
\omega_{\mathrm{CS},\kappa,\mathrm{loc}}$, with at least one compactly
supported argument.

\begin{proposition}[Comparison with abelian Chern--Simons theory]
\label{prop:bf-cs-comparison}
For every nonzero complex $\kappa$, $F_\kappa$ is an isomorphism
of linear BV theories with chiral boundary condition.  It identifies
the local BV Hamiltonian densities and the integrated BF and
level-$\kappa$ abelian Chern--Simons actions on compactly supported
fields satisfying that boundary condition.
This field isomorphism has no specialization at $\kappa=0$.
\end{proposition}
\begin{proof}
The inverse separates the holomorphic degrees and multiplies the
holomorphic-one-form summand by $\kappa$.  Furthermore,
\[
 dF_\kappa(\alpha,\beta)
 =d'\alpha+\partial\alpha+\kappa^{-1}d'\beta
 =F_\kappa Q_\kappa(\alpha,\beta).
\]
The product of two $\alpha$ summands has no three-form component.
The product of two $\beta$ summands has holomorphic degree two and
vanishes.  The two cross terms consequently give
$F_\kappa^*\omega_{\mathrm{CS},\kappa,\mathrm{loc}}
=\omega_{\mathrm{loc}}$ pointwise.
Both $F_\kappa$ and its inverse preserve compact support, so the same
identity holds for the integrated pairings on their stated domains.
The local Hamiltonian densities, the integrated actions on compactly
supported fields, and the boundary condition are preserved.
The inverse scalar $\kappa^{-1}$ in $F_\kappa$ excludes the zero fiber.
\end{proof}

At zero level the field pairing remains the same nondegenerate cross
pairing.  In particular, $b$ remains a field, rather than becoming
$\kappa$ times a fixed field.  Its gauge transformation vanishes,
while the gauge transformation of $a$ remains $d'c$.

This difference can be measured before quantization.  On a product
of a complex disk $D$ and an interior real interval, the local
cohomology of the field complex is computed by
\[
 \mathcal O(D)\xrightarrow{\ \kappa\partial\ }
 \Omega^1_{\mathrm{hol}}(D),
\]
placed in degrees $-1,0$.  The Dolbeault and real Poincare homotopies
remove the antiholomorphic and real directions.  At nonzero $\kappa$,
holomorphic integration on $D$ makes the arrow surjective and its
kernel consists of constants.  At zero it is the zero map.  Thus
\[
\begin{array}{c|cc}
 &H^{-1}(E_\kappa)&H^0(E_\kappa)\\\hline
 \kappa\ne0&\C&0\\
 \kappa=0&\mathcal O(D)&\Omega^1_{\mathrm{hol}}(D).
\end{array}
\]
These are cohomology groups of the field complex.  The negative-degree
summand records gauge data.  The zero fiber therefore has more local
field cohomology than the nonzero Chern--Simons fibers.

\section{A Green operator at zero level}

The zero-level kinetic complex has a concrete propagator operator in
the interior.  Give $\C\times\R$ its flat Euclidean metric and put
\[
 Q^{\mathrm{GF}}=-4\partial_z\iota_{d\bar z}
                         -\partial_t\iota_{dt}.
\]
The contraction symbols satisfy
$\iota_{d\bar z}(d\bar z)=\iota_{dt}(dt)=1$ and kill the other
basis one-forms.  Then
\[
 Q_0Q^{\mathrm{GF}}+Q^{\mathrm{GF}}Q_0
 =-4\partial_z\partial_{\bar z}-\partial_t^2
 =-\partial_x^2-\partial_y^2-\partial_t^2=:\Delta.
\]
The scalar Green kernel is
$G(X,Y)=(4\pi|X-Y|)^{-1}$, with $\Delta_XG=\delta_Y$.
Consequently $H_0=Q^{\mathrm{GF}}G$ satisfies
\[
 Q_0H_0+H_0Q_0=1
\]
on compactly supported test forms, with output interpreted as smooth
forms with the corresponding decay.  Its kernel is
\[
 H_0(X,Y)=
 \frac{\overline{z-w}}{2\pi|X-Y|^3}\iota_{d\bar z}
 +\frac{t-s}{4\pi|X-Y|^3}\iota_{dt},
 \quad X=(z,t),\quad Y=(w,s).
\]
The distribution identity follows from the outward flux of
$-\nabla(4\pi r)^{-1}$ through a sphere, which is one.
Differentiating $G$ gives the two displayed coefficients.
No inverse power of $\kappa$ occurs.  Composing with the inverse
cross pairing gives the BV contraction kernel, including the constant
$2\pi i$ from the normalization of $\omega$.

This interior operator is not a claim about a boundary Green problem.
The boundary comparison below uses compactly supported normal
homotopies which preserve the chiral boundary condition.  Those
homotopies require neither a choice of boundary Green kernel nor a
limit of $\kappa^{-1}$.

\section{The boundary complex and its quantum cocycle}

Write $\mathfrak j(U)=\Omega_c^{0,\bullet}(U)$ for an open
$U\subseteq\Sigma$.  A generator $u\in\Omega_c^{0,q}(U)$ has degree
$q-1$ in the boundary observable complex.  The boundary complement
is the antiholomorphic summand, while the component of the boundary
differential from this complement to $L$ is $\kappa\partial$.
Let $\hbar$ be a central polynomial parameter of cohomological degree
zero.  The boundary observables are
\[
 \mathcal B_\kappa(U)=
 \left(\operatorname{Sym}(\mathfrak j(U)[1])[\hbar],
             \bar\partial+\hbar\Delta_{\mu_\kappa}\right),
 \qquad
 \mu_\kappa(u,v)=-\kappa\int_U^{\mathrm{res}}u\wedge\partial v.
\]
The pairing is zero unless the antiholomorphic degrees sum to one.
As a contraction of observable generators, it has degree one.
The operator $\Delta_{\mu_\kappa}$ vanishes on constants and linear
terms; on a monomial it sums contractions of unordered pairs, with
the Koszul sign for bringing that pair together.  This specifies all
its values.  The differential on a shifted generator is the displayed
$\bar\partial$, transported without an additional minus sign.

The completion in each symmetric power has a concrete meaning.
For the $r$th tensor power, take compactly supported smooth sections
on $U^r$ of the external tensor product of the field bundles; then
pass to coinvariants for the graded action of the symmetric group.
On every fixed compact support use the seminorms of all derivatives,
and take the direct limit over compact supports.  This is the smooth
kernel, or completed bornological tensor, convention.  The symmetric
algebra is the direct sum of these powers.  Thus polynomial degree
is finite; there is no completion in the number of fields.

\begin{lemma}[Boundary quantum differential]
\label{lem:bf-boundary-differential}
The square of $\bar\partial+\hbar\Delta_{\mu_\kappa}$ vanishes
for every $\kappa$.  At $\kappa=0$, the quantum boundary complex is
exactly the commutative Dolbeault algebra
$\operatorname{Sym}(\mathfrak j(U)[1])[\hbar]$.
\end{lemma}
\begin{proof}
For $u$ of antiholomorphic degree $q$, the compatibility equation for
one contraction is
\[
 \mu_\kappa(\bar\partial u,v)
 +(-1)^{q-1}\mu_\kappa(u,\bar\partial v)
 =-\kappa\int_U^{\mathrm{res}}\bar\partial(u\wedge\partial v)=0.
\]
Compact support removes the boundary integral.  It follows that
$\bar\partial\Delta_{\mu_\kappa}+
\Delta_{\mu_\kappa}\bar\partial=0$ on every symmetric power.
In the square of $\Delta_{\mu_\kappa}$, overlapping pairs give zero
because a contracted factor has disappeared.  Disjoint pairs occur
in both orders with opposite signs, since a contraction has odd
degree.  Thus $\Delta_{\mu_\kappa}^2=0$.
Finally, $\mu_0=0$ by its formula.
\end{proof}

The coefficient $\kappa$ comes from the differential
$\kappa\partial$, not from a degeneration of the field pairing.
The next calculation derives this cocycle directly from bulk fields.

\section{A compact collar map for quantum observables}

Fix a smooth compactly supported function $\rho$ on $\R_{\ge0}$
with integral one, and put
$R(t)=\int_0^t\rho(s)ds$.  Thus $R=1$ for sufficiently large $t$.
Let $E_{\kappa,L,c}(U\times\R_{\ge0})$ denote compactly supported
fields satisfying $\iota^*\alpha=0$.
For $u\in\Omega_c^{0,q}(U)$ define
\[
 I_\kappa(u)=u\wedge\rho(t)dt
                +(-1)^q(R(t)-1)\kappa\partial u.
\]
This is a degree-zero map from the unshifted complex
$(\mathfrak j(U),\bar\partial)$ to
$(E_{\kappa,L,c}(U\times\R_{\ge0}),Q_\kappa)$.
The first term is in the $\alpha$ summand and the second is in the
$\beta$ summand.  Both have compact support.  The first has zero
boundary pullback, and the second lies in $L$.

The bulk observable generators are $sx$, for $x\in E_{\kappa,L,c}$,
with $|sx|=|x|-1$.  Define their contraction by
\[
 B(sx,sy)=(-1)^{|x|}\omega(x,y).
\]
For complementary form degrees this is the normalized integral
of the two cross terms.  Let $\Delta_B$ sum pair contractions as above.
The bulk--boundary quantum complex is
\[
 \mathcal O_\kappa(U\times\R_{\ge0})
 =\left(\operatorname{Sym}(E_{\kappa,L,c}[1])[\hbar],
                         Q_\kappa+\hbar\Delta_B\right).
\]
The shifted unary operator is $Q_\kappa(sx)=sQ_\kappa x$.
The pairing identity in Proposition~\ref{prop:bf-bv-family} and the
same disjoint-pair cancellation show that this differential squares
to zero.

\begin{theorem}[The quantum boundary map at every level]
\label{thm:bf-quantum-boundary}
The maps $I_\kappa$ induce quasi-isomorphisms
\[
 \operatorname{Sym}(I_\kappa):\mathcal B_\kappa(U)
       \longrightarrow\mathcal O_\kappa(U\times\R_{\ge0})
\]
for every open $U\subseteq\Sigma$ and every complex $\kappa$,
including zero.  They commute with extension by zero and with
factorization products on disjoint opens.
\end{theorem}
\begin{proof}
First check the unary differential.  Differentiating the first term
of $I_\kappa(u)$ gives
$\bar\partial u\wedge\rho dt+\kappa\partial u\wedge\rho dt$.
Differentiating the second gives
\[
 -(-1)^q(R-1)\kappa\partial\bar\partial u
                          -\kappa\partial u\wedge\rho dt.
\]
The last terms cancel, leaving $I_\kappa(\bar\partial u)$.

Projection to the $\alpha$ summand followed by integration along the
normal one-form gives a cochain map $P$, with $PI_\kappa=1$.
To prove that $P$ is a quasi-isomorphism, consider the exact sequence
whose subcomplex is the $\beta$ summand and whose quotient is the
$\alpha$ summand.  The former has absolute normal forms; the latter
has relative normal forms, whose functions vanish at $t=0$.
On the normal complexes the homotopies are
\[
 \begin{split}
 h_{\mathrm{abs}}(gdt)&=-\int_t^\infty g(s)ds,\\
 h_{\mathrm{rel}}(gdt)&=\int_0^t g(s)ds-R(t)\int_0^\infty g(s)ds.
 \end{split}
\]
Both vanish on functions.  The absolute homotopy satisfies $dh+hd=1$.
The relative one satisfies $dh+hd=1-ip$, where
$p(gdt)=\int g$ and $i(1)=\rho dt$ in degree one.
Their values are compactly supported, and the relative value vanishes
at $t=0$.  Direct differentiation proves all these identities.
For a coefficient of tangential form degree $r$, multiply the
normal homotopy by $(-1)^r$; this cancels the tangential cross terms.
These formulas also apply to smooth sections depending on $z$.
The $\beta$ subcomplex is therefore acyclic, and the relative
$\alpha$ quotient retracts onto $\mathfrak j(U)$ after its displayed
shift.  The long exact sequence proves that $P$, hence $I_\kappa$,
is a quasi-isomorphism.

There is also a continuous contraction for the full split complex.
Let $H$ be the two normal homotopies just displayed, and put
$\delta_\kappa=Q_\kappa-Q_0$.  The correction goes from the
$\alpha$ summand to the $\beta$ summand, so
$\delta_\kappa H\delta_\kappa=0$ and $P\delta_\kappa=0$.
Consequently
\[
 H_\kappa=H-H\delta_\kappa H,\qquad
 I_\kappa=I_0-H\delta_\kappa I_0,\qquad
 Q_\kappa H_\kappa+H_\kappa Q_\kappa=1-I_\kappa P.
\]
To check the last identity, substitute
$Q_0H+HQ_0=1-I_0P$ and
$Q_0\delta_\kappa=-\delta_\kappa Q_0$.
The linear correction terms cancel except for
$H\delta_\kappa I_0P$; the quadratic correction terms vanish.
The formula for $I_\kappa$ is the displayed collar map, since the
absolute homotopy of $\rho dt$ is $R-1$.
This construction uses no division by $\kappa$.

It remains to check the quantum contraction, rather than only the
unary differential.  For antiholomorphic degrees $q+r=1$, expand the
two cross terms in $I_\kappa(u)\wedge I_\kappa(v)$.
Moving $dt$ to the right, and integrating by parts on $U$, gives
\[
 B(sI_\kappa(u),sI_\kappa(v))
 =-2\kappa\left(\int_0^\infty\rho(t)(1-R(t))dt\right)
                          \int_U^{\mathrm{res}}u\wedge\partial v.
\]
The minus sign is the fiber-integration sign fixed above;
$B$ itself integrates the two cross terms with positive sign.
For the second term use
$\int_U\partial u\wedge v=(-1)^{q+1}\int_Uu\wedge\partial v$.
The normal integral is
\[
 \int_0^\infty\rho(1-R)dt
  =\left[R-\frac{R^2}{2}\right]_0^\infty=\frac12.
\]
Consequently $B(I_\kappa(u),I_\kappa(v))=\mu_\kappa(u,v)$, with
suspensions understood.  This proves compatibility on quadratic
monomials.  The pair-contraction definitions prove it on every
symmetric power.

At the classical level, the preceding continuous linear homotopies
extend to $r$ tensor factors by the explicit map
\[
 \mathscr H_r=\sum_{j=1}^r
 (I_\kappa P)^{\otimes(j-1)}\otimes H_\kappa
                       \otimes1^{\otimes(r-j)}.
\]
Use the Koszul sign when the degree-minus-one map $H_\kappa$ passes
the preceding factors.  In $d\mathscr H_r+\mathscr H_rd$, the terms
$1-I_\kappa P$ telescope, giving $1-(I_\kappa P)^{\otimes r}$.
Averaging
over the finite symmetric group gives the result on symmetric powers;
characteristic zero permits this averaging.  The integrations preserve
compact supports in the product variables and extend to the completed
bornological tensors.  Thus $\operatorname{Sym}(I_\kappa)$ is a
classical quasi-isomorphism.

Finally filter both polynomial quantum complexes by the sum of the
power of $\hbar$ and the number of fields.  The unary differential
preserves that filtration degree.  A quantum contraction reduces it
by one.  The associated graded map is the classical
quasi-isomorphism tensored with $\C[\hbar]$.  The filtration starts
in degree zero, and every element has finite filtration degree.
Induction on the finite filtration, followed by the exact direct
limit, proves the quantum quasi-isomorphism.
All formulas use the same $\rho$ and integrate only the normal
variable.  They therefore commute with extension by zero and with
products on disjoint boundary opens.
\end{proof}

The observables also satisfy the gluing property needed for a
factorization algebra.  The relevant covers $\{U_i\}$ of $U$ are
those for which every finite subset of $U$ lies in some $U_i$.
For every $r$, the sets $U_i^r$ then cover $U^r$.
At a fixed field degree the augmented Cech complex of compactly
supported smooth kernels for this cover is exact: a locally finite
partition of unity gives its contraction, and only finitely many
terms meet a given compact support.  Multiplication by the partition
functions preserves the linear boundary condition for bulk fields.
For fixed $r$ the form degrees are bounded, so the finite filtration
by form degree upgrades this Cech exactness to descent for the unary
differential.  Passing to symmetric coinvariants preserves exactness
because averaging over the finite symmetric group is available.
The polynomial filtration used above then proves quantum descent.
Here factorization algebra means precisely these gluing complexes
together with the compatible products for disjoint opens.
Thus the theorem is an equivalence of quantum factorization algebras
on the boundary.  At $\kappa=0$ its image contains only $\alpha$
fields, so its pair contraction vanishes directly.

\section{A compact calculation of the current level}

The quantum differential determines how a product changes when its
insertions are moved.  A compact two-current calculation fixes the
sign before any mode presentation is chosen.
For a point $a$ in a disk, choose a cutoff $\chi_a$ supported in a
small disk about $a$ and equal to one near its center.  The form
\[
 u_a=-\frac{\bar\partial\chi_a}{z-a}
\]
is smooth and compactly supported.  Stokes' theorem on an annulus
about $a$ gives, for every holomorphic function $F$ on the disk,
\[
 \int^{\mathrm{res}}u_a\wedge F(z)dz=F(a).
\]
Thus $u_a$ is a smooth representative of evaluation at $a$.

Choose distinct points $a,b$, disjoint small disks about them, and
cutoffs $\chi_a,\chi_b$ as above.  Choose a common radial cutoff
$\chi(z)=\psi(|z|^2)$, equal to one on both small disks and supported
in a larger disk centered at zero.  Put
\[
 e_a=-\frac{\bar\partial\chi}{z-a},\quad
 e_b=-\frac{\bar\partial\chi}{z-b},\qquad
 f=\frac{\chi-\chi_a}{z-a},\quad
 g=\frac{\chi-\chi_b}{z-b}.
\]
The functions $f,g$ are smooth and compactly supported: the apparent
poles are removable.  Direct differentiation gives
$\bar\partial f=u_a-e_a$ and $\bar\partial g=u_b-e_b$.
On the support of $u_b$, $f(z)=(z-a)^{-1}$.  Integration by parts
and the evaluation formula give
\[
 \mu_\kappa(f,u_b)
 =-\kappa\int^{\mathrm{res}}u_b\wedge\partial f
 =-\kappa f'(b)=\frac{\kappa}{(a-b)^2}.
\]
The other replacement has zero contraction.  On the common cutoff
annulus, $g=\chi/(z-b)$ and
\[
 e_a\wedge\partial g
 =-\frac{\bar\partial\chi\wedge\partial\chi}{(z-a)(z-b)}
  +\frac{\chi\bar\partial\chi\wedge dz}{(z-a)(z-b)^2}.
\]
Relative to $d\bar z\wedge dz$, its coefficient is
\[
 -\frac{(\psi')^2|z|^2}{(z-a)(z-b)}
 +\frac{\psi\psi' z}{(z-a)(z-b)^2}.
\]
Both Laurent expansions have strictly negative angular Fourier powers,
beginning with $z^{-2}$.  Their angular integrals vanish, so
$\mu_\kappa(e_a,g)=0$.

The function generators have degree $-1$ and the form generators have
degree zero.  With the specified positive unary differential,
$D=\bar\partial+\hbar\Delta_{\mu_\kappa}$, the Leibniz signs and
these contractions give the compact chain relation
\[
 D(fu_b+e_ag)=u_au_b-e_ae_b+
                     \frac{\hbar\kappa}{(a-b)^2}.
\]
Consequently the actual separated product satisfies
\[
 [u_au_b]=[e_ae_b]-\frac{\hbar\kappa}{(a-b)^2}\mathbf1.
\]
The radial term has the regular expansion
$e_a=\sum_{n\ge0}a^nu_n$, where
$u_n=-\bar\partial\chi/z^{n+1}$, and similarly for $e_b$.
These series converge in every smooth seminorm on the outer annulus.
They cannot contribute a diagonal pole.  The current determined by
the fixed state $q_1=[u_0]$ therefore has parameter $-\kappa$:
\[
 J(a)J(b)\sim-\frac{\hbar\kappa}{(a-b)^2}.
\]
The remaining problem is to obtain every field and regular product
from the supported complex, keeping this normalization.

A direct annular evaluation of the local cocycle has a different
role.  For holomorphic Laurent functions $F,G$ and a cutoff equal to
one at the inner boundary and zero at the outer boundary,
\[
 \mu_\kappa(F,\bar\partial\chi\,G)
 =-\kappa\int^{\mathrm{res}}\bar\partial\chi\wedge G\,dF
 =\kappa\operatorname{Res}G\,dF.
\]
Only one factor needs compact support in this integral.  The minus
sign in the separated product comes from moving the quantum
contraction to the other side of an exact chain relation.  Thus the
annular scalar evaluation by itself does not identify the extracted
mode action.  Multiplying all currents by $i$ changes their parameter,
but also changes the prescribed state and its jet representatives.

\section{Cauchy contraction on smooth observables}

The quantum differential can be removed by an explicit linear map,
provided its effect on products is retained.  For
$u\in\Omega_c^{0,1}(U)$, extend $u$ by zero to $\C$ and define its
Cauchy transform by
\[
 (\mathscr C u)(z)=\frac1{2\pi i}\int_{\C}
                       \frac{u(w)\wedge dw}{z-w}.
\]
The kernel $1/(z-w)$ is locally integrable.  Stokes' theorem on a
small circle gives
$\partial_{\bar z}(1/(\pi z))=\delta_0$; the normalization
$d\bar w\wedge dw=2i\,dA_w$ then gives
\[
 \bar\partial\mathscr C u=u,\qquad
 \mathscr C\bar\partial f=f\quad(f\in C_c^\infty(\C)).
\]
For the second identity, the first implies that the difference is
holomorphic.  Integration by parts shows that all moments of
$\bar\partial f$ vanish, so expansion of the Cauchy kernel outside
a disk containing the support makes $\mathscr C\bar\partial f$ compactly supported.
The difference is therefore a compactly supported holomorphic
function and is zero.  The output $\mathscr C u$ is generally not
compactly supported; it is used below only inside a scalar integral.

On two degree-zero observable generators, define
\[
 P_\kappa(u,v)=-\kappa\int^{\mathrm{res}}
                             (\mathscr C u)\,\partial v.
\]
Define $P_\kappa$ to be zero on all other generator-degree pairs.
The derivative of $v$ has compact support, so the displayed integral
exists.  If $u=a(w)d\bar w$ and $v=b(z)d\bar z$, it is
\[
 P_\kappa(u,v)=\frac{\kappa}{\pi^2}
     \int_{\C^2}\frac{a(w)\,\partial_z b(z)}{z-w}\,dA_zdA_w.
\]
This formula also defines a pair contraction on an arbitrary compactly
supported smooth kernel, by differentiating in the paired $z$
variable and integrating both paired variables.  The output is a
smooth kernel on the remaining variables, supported in the compact
projection of the input support.  For each fixed support $K$ and
multi-index $\alpha$ in the remaining variables,
\[
 \|\partial^\alpha\Delta_P F\|_\infty
       \le C_{K,\alpha}\|F\|_{C^{|\alpha|+1}}.
\]
Indeed, the integral of $|z-w|^{-1}$ on a fixed compact pair of planar
variables is finite.  Derivatives in the remaining variables pass
under that integral.  Thus the contraction is continuous in the
smooth-kernel convention already specified; no singular observable
or extra completion is introduced.

The pairing $P_\kappa$ is symmetric.  After interchanging $z,w$, its
antisymmetric part is proportional to
\[
 \int_{\C^2}
 \frac{(\partial_z+\partial_w)(a(w)b(z))}{z-w}\,dA_zdA_w.
\]
In difference and translation coordinates, the derivative in the
numerator differentiates only the translation coordinate.  Its
integral is zero by compact support, without differentiating the
singular difference kernel.

Let $K=\Delta_{P_\kappa}$ sum unordered pair contractions, with the
shifted Koszul signs, and write $d=\bar\partial$.  The operator $K$
has degree zero and lowers field number by two.  Its commutators are
\[
 [d,K]=-\Delta_{\mu_\kappa},\qquad
 [K,\Delta_{\mu_\kappa}]=0.
\]
To verify the first equation, it suffices first to use a function
and a form generator:
\[
 [d,K](fu)=-P_\kappa(\bar\partial f,u)
                       =-\mu_\kappa(f,u).
\]
The other quadratic generator degrees give zero.  Both sides are
constant pair-contraction operators, so the identity holds on all
monomials.  For the second equation, overlapping contractions vanish
because the shared factor disappears; disjoint contractions commute
because $K$ is even.  The same identities hold on smooth kernels by
continuity.  One may verify density of decomposable kernels on product
boxes by Fourier expansion after inserting compact cutoffs: rapid
decay of Fourier coefficients gives convergence in every derivative
seminorm.  A finite partition of unity on the input support gives
the general case.

\begin{proposition}[Finite quantum conjugation]
\label{prop:bf-cauchy-conjugation}
For every open $U\subseteq\C$ and every $\kappa\in\C$, the map
\[
 N_U=e^{-\hbar K}:
 (\mathcal B_\kappa(U),D)
       \longrightarrow
 (\operatorname{Sym}(\Omega_c^{0,\bullet}(U)[1])[\hbar],d)
\]
is a chain isomorphism, with inverse $e^{\hbar K}$.  It commutes
with extension by zero, translations and rotations.
\end{proposition}
\begin{proof}
The two commutators give
\[
 e^{\hbar K}d\,e^{-\hbar K}
   =d+\hbar[K,d]=d+\hbar\Delta_{\mu_\kappa}=D.
\]
On field number $r$, $K^j=0$ for $2j>r$, so all exponentials terminate.
They preserve the polynomial complexes and their filtration by
field number plus $\hbar$-power.  The global Cauchy formula uses the
same zero extension for every open set.  Its kernel and measure are
translation invariant and transform compatibly under rotations.
These facts prove the stated naturality.
\end{proof}

The map $N_U$ is not asserted to preserve multiplication at overlapping
supports.  Its cross contractions carry the quantum products.

\section{The vacuum states and their smooth representatives}

Let $D_R=\{|z|<R\}$ and choose a radial cutoff $\chi$ equal to one
near zero and zero near the boundary.  For $n\ge0$ put
\[
 u_n=-\frac{\bar\partial\chi}{z^{n+1}}.
\]
The support avoids zero.  Stokes' theorem gives its divided-jet moments:
\[
 \int_{D_R}^{\mathrm{res}}u_n\wedge F(z)dz
       =\frac{F^{(n)}(0)}{n!},\qquad
 M_j(u_n)=\delta_{jn},\quad
 M_j(u)=\int_{D_R}^{\mathrm{res}}u\wedge z^jdz.
\]
A compactly supported holomorphic function vanishes.  A compact
$(0,1)$-form with every holomorphic moment zero has a compact primitive:
its Cauchy transform vanishes outside a disk containing its support,
by expansion in $w/z$.  Conversely, integration by parts kills the
moments of $\bar\partial f$ for compact $f$.

The required statement on completed tensors follows from an explicit
estimate.  If the input is supported in $|z|\le r_0<R$, choose the
cutoff annulus with inner radius $\sigma>r_0$.  Then
\[
 |M_n(u)|\le C r_0^n,\qquad
 \|u_n\|_{C^d}\le C_d(n+1)^{d+1}\sigma^{-n-1}.
\]
Thus the maps
\[
 \Pi u=\sum_{n\ge0}M_n(u)u_n,\qquad
 h_Du=\mathscr C(u-\Pi u)
\]
are well defined on this fixed support space, with convergence in
all derivative seminorms and compact output.  On functions set
$\Pi=h_D=0$.  The moments and the Cauchy identities prove
\[
 dh_D+h_Dd=1-\Pi.
\]
For a compact smooth kernel on $D_R^r$, choose a common inner support
radius and a larger cutoff annulus in each variable.  The moment
bounds multiply, so the resulting multiple series converge in every
mixed derivative seminorm.  With shifted tensor signs, the map
\[
 H_r=\sum_{j=1}^r
       \Pi^{\otimes(j-1)}\otimes h_D\otimes1^{\otimes(r-j)}
\]
satisfies $dH_r+H_rd=1-\Pi^{\otimes r}$ by telescoping.
Each variable is treated once.  No single cutoff operator on all
compact supports is required.  Averaging over the finite symmetric
group proves the same assertion for the completed symmetric powers.

The tensor cohomology is consequently in top antiholomorphic degree
$r$, hence in shifted cohomological degree zero, and is determined
by its multivariable moments.  This computes the odd generators by
their differential; it does not discard them.

Rotations act on the representatives by
\[
 R_\theta^*u_n=e^{-i(n+1)\theta}u_n.
\]
Define the state weight as minus this pullback rotation weight.
The radial homotopies commute with rotations, including simultaneous
rotation of all variables.  At state weight $N$, a surviving tensor
must satisfy
\[
 \sum_{i=1}^r(n_i+1)=N.
\]
There are finitely many such tuples, and $r\le N$.  Thus the argument
applies to a fixed diagonal weight without assuming a finite Fourier
expansion in the individual variables.

Rotation invariance also gives $P_\kappa(u_n,u_m)=0$ for all $n,m$:
the two inputs have strictly positive total state weight, while a
scalar has weight zero.  The map $N_{D_R}$ therefore fixes every
monomial in these radial representatives.  Combining its chain
isomorphism with the moment calculation gives
\[
 V_\kappa=\C[q_1,q_2,\ldots][\hbar],\qquad
 q_{n+1}\longmapsto[u_n],\quad
 |q_m|=0,\quad\mathrm{wt}(q_m)=m.
\]
More precisely, $V_\kappa$ is the direct sum of finite-weight
cohomology spaces of $\mathcal B_\kappa(D_R)$.  Each weight space is
finite free over $\C[\hbar]$, and there is no odd cohomology or
$\hbar$-torsion.  The full disk cohomology also contains infinite
moment sequences and is not identified with this direct sum.

Polynomial multiplication specifies the normal state basis
$\prod q_{m_i}\mapsto[\prod u_{m_i-1}]$.  It does not assert that
multiplication at overlapping supports descends to quantum
cohomology.  Changing the radial cutoff or enlarging the centered
disk preserves this basis: applying $N$ reduces the comparison to
the equality of its multivariable moments.

\section{Transport of every separated product}

For disjoint opens $U_1,\ldots,U_s\subseteq U$, the original
factorization product is extension by zero followed by multiplication
$m$.  Its mixed $\mu_\kappa$ contractions vanish because the supports
are disjoint.  Under Cauchy conjugation its exact formula is
\[
 N_Um(N_{U_1}^{-1}\otimes\cdots\otimes N_{U_s}^{-1})
   =m\exp\left(-\hbar\sum_{i<j}\Delta_{P_\kappa}^{ij}\right),
\]
where $\Delta_{P_\kappa}^{ij}$ contracts a pair with one factor in
input $i$ and one in input $j$.  To prove the formula, decompose $K$
on a product into contractions internal to each input and contractions
between inputs.  All these even operators commute.  The internal
exponentials cancel $N_{U_i}^{-1}$, leaving exactly the displayed
cross contractions.

The exponential terminates on every input.  Its terms are matchings
of the individual field factors, with edges joining different input
groups.  A matching with $j$ edges appears $j!$ times, one for each
ordering of its disjoint contractions; the exponential coefficient
$1/j!$ cancels that multiplicity.  The proof extends to all compact
smooth kernels by the continuity estimate above.  It commutes with
nesting, either by conjugating the original factorization maps or
by observing that each nested expression contains each cross pair
exactly once.

For a small radial cutoff at $a$, write
\[
 u_{a,n}=-\frac{\bar\partial\chi_a}{(z-a)^{n+1}},\qquad n\ge0.
\]
Its moments are divided derivatives at $a$.  Outside the supporting disk,
\[
 (\mathscr C u_{a,n})(z)=\frac1{(z-a)^{n+1}}.
\]
This follows by applying the moment formula to the holomorphic
function $w\mapsto(z-w)^{-1}$.  For disjoint supporting disks at
$a,b$, integration by parts therefore gives
\[
 \begin{split}
 P_\kappa(u_{a,n},u_{b,m})
 &=-\frac\kappa{m!}\left.
       \partial_z^{m+1}(z-a)^{-n-1}\right|_{z=b}\\
 &=\frac\kappa{n!m!}
       \partial_a^n\partial_b^m(a-b)^{-2}.
 \end{split}
\]
Put $\lambda=-\hbar\kappa$.  The actual cross contraction of divided
jets of orders $n,m$ is
\[
 C_{n,m}(a,b)=
 \lambda\frac{(-1)^n(n+m+1)!}{n!m!}(a-b)^{-n-m-2}.
\]
In particular, $C_{0,0}(a,b)=\lambda(a-b)^{-2}$, in agreement with
the compact two-current calculation.

Holomorphic motion of the remaining state is equally explicit.  A
common outer radial cutoff gives
\[
 -\frac{\bar\partial\chi}{(z-a)^{n+1}}
  =\sum_{k\ge0}\binom{n+k}{k}a^ku_{n+k}.
\]
The difference from $u_{a,n}$ has zero moments and a compact primitive.
For products of co-centered representatives, all internal $P$ terms
vanish by rotation about that center.  The same comparison therefore
holds for each polynomial state.  It gives the translation derivation
\[
 Tq_m=mq_{m+1},\qquad T1=0.
\]
The moving state corresponding to $q_m$ is $e^{aT}q_m$.

For arbitrary polynomial states inserted at distinct points, label
each of their jet factors by $\alpha$, with jet order $r_\alpha$ and
insertion group $i(\alpha)$.  The actual product, expressed by $N$
and the moment coordinates, is
\[
 \sum_M
  \left(\prod_{\{\alpha,\beta\}\in M}
       C_{r_\alpha,r_\beta}(a_{i(\alpha)},a_{i(\beta)})\right)
  \left(\prod_{\gamma\ \mathrm{unpaired}}
                e^{a_{i(\gamma)}T}q_{r_\gamma+1}\right).
\]
Here $M$ ranges over matchings joining different insertion groups.
There are finitely many matchings, including the empty matching.
The empty matching retains the regular terms.
For fixed numerical insertion points, the unpaired series have
smooth representatives on a common outer annulus and converge in
all derivative seminorms, by the geometric moment bounds.  Thus the
formula is an equality in full disk cohomology.  Taking weight
coefficients and radial Laurent coefficients of this equality gives
polynomial states.  This is the extraction from the smooth carrier.

\section{Fields, modes and regular products}

Applying the separated-product formula to one current and a monomial
at zero gives the operator-valued formal distribution
\[
 J(z)=\sum_{m\ge1}q_m z^{m-1}
           +\lambda\sum_{m\ge1}m\partial_{q_m}z^{-m-1}.
\]
The first sum is the creation part $J_-(z)$, acting by multiplication;
the second is the annihilation part $J_+(z)$, acting by derivations.
A mode $a_{(n)}$ is the coefficient of $z^{-n-1}$ in its field.
For this current,
\[
 J_{-m}=q_m,\qquad
 J_m=-\hbar\kappa m\partial_{q_m}\quad(m>0),\qquad J_0=0.
\]
These operators were obtained from compact insertions and their
quantum differential with the fixed identification $q_{n+1}=[u_n]$.

Write $J^{[r]}(z)=\partial_z^rJ(z)/r!$ for its divided derivatives.
For a monomial state $a=\prod_{i=1}^r q_{m_i}$, its field is
\[
 Y(a,z)=:\prod_{i=1}^rJ^{[m_i-1]}(z):.
\]
Normal ordering here has an explicit definition.  Split each divided
current into creation and annihilation parts, expand the finite
product, and put every multiplication operator to the left of every
derivation.  Both types commute within themselves, so the answer is
independent of the ordering of the factors of $a$.  Expanding those
derivations against any polynomial state gives precisely the preceding
matching formula with the second insertion at zero.  Thus $Y$ is the
extracted state-field operation, including its regular coefficients.

For example, the complete two-jet product is
\[
 \begin{split}
 Y(q_m,z)q_n
 &=\sum_{k\ge0}\binom{m+k-1}{k}q_{m+k}q_nz^k\\
 &\quad+\lambda\frac{(-1)^{m-1}(m+n-1)!}
                 {(m-1)!(n-1)!}\,z^{-m-n}.
 \end{split}
\]
For a quadratic state, the regular coefficient contains a contraction
with a translated unpaired current:
\[
 (q_1^2)_{(-1)}q_1=q_1^3+2\lambda q_3
                  =q_1^3-2\hbar\kappa q_3.
\]
Indeed,
$Y(q_1^2,z)q_1=J_-(z)^2q_1+2\lambda z^{-2}J_-(z)$.
Since $q_{1(-1)}q_1=q_1^2$ and $q_{1(-1)}q_1^2=q_1^3$,
this also gives the regular-product associator
\[
 (q_{1(-1)}q_1)_{(-1)}q_1
       -q_{1(-1)}(q_{1(-1)}q_1)=2\lambda q_3.
\]
Polynomial multiplication in the normal state basis and the regular
vertex product are therefore different operations at nonzero level.

Every field is lower truncated on every polynomial state: only
finitely many annihilation modes can differentiate that state.
At any fixed remaining nonnegative power of $z$, its creation factors
have only finitely many index choices.  All field coefficients are
therefore well-defined polynomial endomorphisms.

\section{All operator products and the vertex identities}

Direct differentiation on the polynomial states gives
\[
 [J_m,J_n]=\lambda m\delta_{m+n,0},\qquad
 [J_+(z),J_-(w)]=\iota_{z,w}\frac\lambda{(z-w)^2},
\]
where $\iota_{z,w}$ means expansion in $w/z$.
Commuting annihilation factors past creation factors in the normal
fields proves their complete Wick formula.  If the factors of $a$
have jet orders $r_\alpha$ and those of $b$ have jet orders $s_\beta$,
then
\[
 Y(a,z)Y(b,w)=
 \sum_M\iota_{z,w}\left(
          \prod_{\{\alpha,\beta\}\in M}
                        C_{r_\alpha,s_\beta}(z,w)\right)
       :\!\prod_{\alpha\ \mathrm{unpaired}}J^{[r_\alpha]}(z)
          \prod_{\beta\ \mathrm{unpaired}}J^{[s_\beta]}(w)\!:\ .
\]
Each matching joins a factor of $a$ to a factor of $b$; unpaired
fields keep their original positions and divided derivatives.
Every commutation removes a paired creation and annihilation factor,
so induction on the number of factors proves the formula and its
matching multiplicities.  It also proves the analogous formula for
any number of insertion groups.  For instance,
\[
 Y(q_1^2,z)Y(q_1^2,w)
 =:J(z)^2J(w)^2:
   +\frac{4\lambda}{(z-w)^2}:J(z)J(w):
   +\frac{2\lambda^2}{(z-w)^4},
\]
with each rational coefficient expanded by $\iota_{z,w}$.
Taylor expansion of the unpaired factors gives every regular as well
as every singular operator-product coefficient.

Set $\mathbf1=1$.  The annihilation operators kill the vacuum, and
the creation part of $J^{[m-1]}(z)$ is $e^{zT}q_m$.  Hence
\[
 Y(\mathbf1,z)=\mathrm{id},\qquad
 Y(a,z)\mathbf1=e^{zT}a,
 \qquad T\mathbf1=0.
\]
The translation identities follow from
\[
 [T,q_m]=m q_{m+1},\qquad
 [T,\partial_{q_m}]=-(m-1)\partial_{q_{m-1}},
\]
where the second expression is zero at $m=1$.  They give
\[
 [T,Y(a,z)]=\partial_zY(a,z)=Y(Ta,z).
\]

Locality also follows from the same matching formula.  The reverse
ordered product has the same normally ordered unpaired factors and
the reverse expansion of the same rational contractions.  If $a,b$
have $r,s$ factors, respectively, put
\[
 L=\sum_\alpha r_\alpha+\sum_\beta s_\beta+2\min(r,s).
\]
Every matching denominator divides $(z-w)^L$.  After multiplication
by this factor the two radial expansions agree, proving
$(z-w)^L[Y(a,z),Y(b,w)]=0$.

The iterate identity requires compatibility of the regular terms as
well.  Introduce a formal displacement $x$ and put $z=w+x$.
Each unpaired jet expands as
\[
 J^{[r]}(w+x)=\sum_{k\ge0}\binom{r+k}{k}x^kJ^{[r+k]}(w).
\]
These are exactly the translated jet states in $Y(a,x)b$.
For any third polynomial state $c$, matchings between the three groups
$a,b,c$, at $z,w,0$, define one expression
\[
 F_{a,b,c}(z,w)\in
 V_\kappa[[z,w]][z^{-1},w^{-1},(z-w)^{-1}]
\]
with finite pole orders.  Its three expansions are
\[
 \begin{aligned}
 \iota_{z,w}F&=Y(a,z)Y(b,w)c,\\
 \iota_{w,z}F&=Y(b,w)Y(a,z)c,\\
 \iota_{w,x}F(w+x,w)&=Y(Y(a,x)b,w)c.
 \end{aligned}
\]
For the last identity, group a matching first by its $a$--$b$ edges
and then by its edges to $c$.  This retains each matching and its
multiplicity.  At fixed output weight, the expression is a rational
function with polynomial numerator in a finite free weight space.
This makes all three expansions legitimate coefficient calculations.

For completeness, Jacobi follows from an elementary formal identity.
Write $\delta(t)=\sum_{n\in\Z}t^n$ and
$\delta(z,w)=z^{-1}\delta(w/z)$.  With the corresponding binomial
expansion orders,
\[
 \begin{split}
 x^{-1}\delta\left(\frac{z-w}{x}\right)
 -x^{-1}\delta\left(\frac{w-z}{-x}\right)
 =w^{-1}\delta\left(\frac{z-x}{w}\right).
 \end{split}
\]
The coefficients of negative powers of $x$ cancel on the left,
because they come from nonnegative powers of $z-w$.  For $k\ge0$,
the coefficient of $x^k$ on each side is
$\partial_w^k\delta(z,w)/k!$; this is the binomial identity
$\binom{-l-1}{k}=(-1)^k\binom{l+k}{k}$.
Multiply the identity by each rational monomial of a fixed-weight
component of $F_{a,b,c}$, using $z=w+x$ in its collision expansion.
The same binomial calculation, including negative exponents, gives
\[
 \begin{split}
 &x^{-1}\delta\left(\frac{z-w}{x}\right)Y(a,z)Y(b,w)
 -x^{-1}\delta\left(\frac{w-z}{-x}\right)Y(b,w)Y(a,z)\\
 &\qquad=w^{-1}\delta\left(\frac{z-x}{w}\right)
                                     Y(Y(a,x)b,w).
 \end{split}
\]
This is the full vertex Jacobi identity.  In modes it reads, for all
integers $m,n,k$,
\[
 \begin{split}
 \sum_{j\ge0}\binom mj(a_{(n+j)}b)_{(m+k-j)}
 =\sum_{j\ge0}(-1)^j\binom nj
 \bigl(a_{(m+n-j)}b_{(k+j)}
       -(-1)^n b_{(n+k-j)}a_{(m+j)}\bigr).
 \end{split}
\]
The sums are locally finite on polynomial states by lower truncation.
All extracted states are even, as the cohomology calculation proved.

\begin{theorem}[The boundary current at every level]
\label{thm:bf-current-all-levels}
For every $\kappa\in\C$, the polynomial finite-weight chiral vacuum
extracted from $\mathcal B_\kappa$ is a vertex algebra over
$\C[\hbar]$.  Its state space, translation and fields are
\[
 V_\kappa=\C[q_1,q_2,\ldots][\hbar],\qquad
 Tq_m=mq_{m+1},\qquad
 Y\left(\prod_iq_{m_i},z\right)=:\prod_iJ^{[m_i-1]}(z):,
\]
where $J_{-m}=q_m$, $J_m=-\hbar\kappa m\partial_{q_m}$ for $m>0$,
and $J_0=0$.  The identification $q_{n+1}=[u_n]$ preserves every
mode, divided jet, regular product and separated factorization
product.  Its Heisenberg level is $-\kappa$ with this fixed current.
\end{theorem}
\begin{proof}
The continuous Cauchy conjugation and compact moment homotopies give
the stated cohomology and its normal state basis.  Transport through
that conjugation gives the actual separated-product matching formula.
Expanding it with one input at zero yields exactly the fields above,
including their regular terms.  The preceding vacuum, translation,
locality and Jacobi calculations prove all vertex identities.
The chain relation and the mode calculation both give level
$-\kappa$ without changing the chosen state representatives.
\end{proof}

\begin{theorem}[The zero-level current]
\label{thm:bf-zero-current}
At $\kappa=0$, the extracted vertex algebra is the commutative
differential algebra $\C[q_1,q_2,\ldots][\hbar]$ with
$Tq_m=mq_{m+1}$ and $Y(a,z)b=(e^{zT}a)b$.  It agrees with the
algebraic zero-level Heisenberg vacuum in every mode and product.
\end{theorem}
\begin{proof}
Both quantum contractions and the annihilation part of $J$ vanish
at zero.  Only the empty matching remains.  Its state-field formula
is precisely $(e^{zT}a)b$.
\end{proof}

\begin{corollary}[Augmentation of the extracted current]
The zero-level polynomial vertex algebra has the augmentation
$q_m\mapsto0$ for all $m\ge1$.  At nonzero level there is no
augmentation to $\C[\hbar]$ as a vertex algebra over $\C[\hbar]$.
\end{corollary}
\begin{proof}
At zero, evaluation at $q_m=0$ preserves multiplication, translation
and the vacuum.  At nonzero level $q_{1(1)}q_1=-\hbar\kappa\mathbf1$.
A vertex map to the scalar vertex algebra would send its left side
to zero and its right side to a nonzero scalar.
\end{proof}

\section{Formal completion and the bulk--boundary comparison}

The preceding constructions are polynomial in $\hbar$ and in field
number.  The Cauchy exponentials terminate on every observable, and
every fixed tuple of polynomial states has finitely many matchings
and finite pole orders.  Smooth moment expansions converge on a
common outer annulus before their finite-weight coefficients are
extracted.

Only after this polynomial extraction take the coefficientwise
$\hbar$-completion
\[
 \widehat V_\kappa
 =\varprojlim_M V_\kappa/\hbar^MV_\kappa
 =\C[q_1,q_2,\ldots]\llbracket\hbar\rrbracket.
\]
Every coefficient of $\hbar$ is a finite polynomial, while its field
number and weight may increase with that coefficient.  The completed
state-field operation has target
\[
 \varprojlim_M
       (V_\kappa/\hbar^MV_\kappa)((z)).
\]
Modulo each $\hbar^M$, the inputs are polynomial states and all the
preceding identities hold.  The operations commute with reduction
in $M$.  Vacuum, translation, Jacobi and locality modulo each
$\hbar^M$ therefore pass to this inverse limit.

There need not be a uniform Laurent pole bound for completed inputs.
At $\kappa\ne0$, take
$a=\sum_{n\ge0}\hbar^nq_{n+1}$ and $b=q_1$.
Its singular terms are
\[
 -\kappa\sum_{n\ge0}(-1)^n(n+1)\hbar^{n+1}z^{-n-2}.
\]
The pole orders increase with the $\hbar$-order.  Thus the inverse
limit above is the correct field target, rather than
$\widehat V_\kappa((z))$ with a single Laurent bound.  This passage
completes the extracted state and field algebra; it makes no claim
that cohomology commutes with completion of the entire factorization
algebra.

The supported quantum collar map of
Theorem~\ref{thm:bf-quantum-boundary} has its original fields,
differential and contraction.  It commutes with translations,
rotations and every separated product.  Hence the current states
in the bulk--boundary pushforward are
\[
 q_m\longmapsto[I_\kappa(u_{m-1})],\qquad
 \prod_iq_{m_i}\longmapsto
                  \left[\prod_i I_\kappa(u_{m_i-1})\right].
\]
Preservation of separated products passes to their moving-point
coefficients, so this comparison preserves every extracted field
and mode.  The nonzero BF--Chern--Simons field change has bulk
parameter $\kappa$, whereas the fixed boundary current has level
$-\kappa$ with the printed positive unary differential.  At zero,
this same supported construction realizes the commutative current
while retaining the additional bulk gauge and field cohomology.

\section{Separated products and topological higher operations}

For distinct holomorphic insertion points, the zero-level Wick
product has no contraction terms.  Its normal-ordering map sends
\[
 x_{i,r}\longmapsto\partial_{z_i}^rJ(z_i),\qquad
 x_{1,0}\star x_{2,0}\star x_{3,0}
       =x_{1,0}x_{2,0}x_{3,0}.
\]
These are equalities on the actual boundary field representation
from Theorem~\ref{thm:bf-zero-current}.

An auxiliary topological interval gives the same higher operations
on these fields as on their algebraic presentation.  To specify the
statement, let $e_0,e_1$ be endpoint cochains and $\eta$ the oriented
edge, with $de_0=-\eta$, $de_1=\eta$.  Their polynomial-form
inclusion and homotopy are
\[
 i(e_0)=1-t,\quad i(e_1)=t,\quad i(\eta)=dt,
 \qquad
 h(gdt)=\int_0^t g-t\int_0^1g.
\]
Projection takes endpoint values and the integral of a one-form.
Direct differentiation gives $dh+hd=1-ip$ and $hi=ph=h^2=0$.
Extend these maps by the identity on boundary field coefficients.

To check all higher identities, let $W$ be the interval cochains
with these field coefficients and let $A$ be their polynomial-form
resolution.  The coefficient algebra is the commutative algebra of
represented zero-level current fields at the chosen separated points,
interpreted coefficientwise in the insertion variables.  Write $sa$
for a suspension, with $|sa|=|a|-1$.  Put
\[
 D_1(sa)=s(da),\qquad D_2(sa,sb)=(-1)^{|a|}s(ab).
\]
Extend these maps to the direct sum of finite tensor words by inserting
them into each consecutive block, with the Koszul sign from the
preceding suspended inputs.  This is a coderivation for the coproduct
which splits words into initial and final segments.  Associativity
and the Leibniz rule say $(D_1+D_2)^2=0$.
Use the same letters $i,p,h$ for the suspended linear maps and define
\[
 F_1=i,\qquad R_n=\sum_{a+b=n}D_2(F_a\otimes F_b),\qquad
 F_n=-hR_n,\qquad b_n=pR_n\quad(n\ge2),
\]
with $b_1=sds^{-1}$.  An $A_\infty$ structure here means that the
resulting degree-one coderivation $b$ on the tensor words in $sW$
satisfies $b^2=0$.  The ordinary operations are fixed by
\[
 b_n(sa_1,\ldots,sa_n)
 =(-1)^{\sum_j(n-j)|a_j|}s m_n(a_1,\ldots,a_n).
\]

Here is a direct verification of the recursion.  The maps $F_n$
extend uniquely to a coalgebra map $F$.  Induct on word length,
assuming $DF=Fb$ and $b^2=0$ in smaller lengths, where $D=D_1+D_2$.
With the length-$n$ terms temporarily zero, the length-$n$ morphism
defect is
\[
 Q_n=R_n-L_n,\qquad
 L_n=\sum_{\substack{r+q+t=n\\2\le q\le n-1}}
 F_{r+1+t}(1^{\otimes r}\otimes b_q\otimes1^{\otimes t}).
\]
The tensor substitutions carry the insertion signs just specified.
Let $O_n$ be the length-$n$ term of $b^2$ with $b_n=0$.
The identity
\[
 D(DF-Fb)+(DF-Fb)b=-Fb^2
\]
gives $\delta Q_n=-iO_n$, where
$\delta Q_n=D_1Q_n+Q_nb_1$ and $b_1$ also denotes its extension to
words.  For $k\ge2$, $pF_k=hF_k=0$, by $ph=h^2=0$.
It follows that $pL_n=hL_n=0$.
Setting $b_n=pQ_n=pR_n$ gives $\delta b_n=-O_n$, which is the
missing length-$n$ identity for $b^2$.
Setting $F_n=-hQ_n=-hR_n$ and using $D_1h+hD_1=1-ip$ gives
\[
 \delta F_n=-Q_n+ipQ_n+h\delta Q_n=ib_n-Q_n,
\]
because $hi=0$.  This is the missing morphism identity.
The induction starts with the cochain maps $i,p$, so it proves every
arity without an infinite sum on any word.  The linear map $i$ is a
quasi-isomorphism by the explicit contraction.

Every field comparison intertwines each multiplication vertex and
each interval homotopy separately.  It therefore intertwines these
operations in every arity, before summing their binary trees.
For example,
\[
 h(tdt)=\frac{t^2-t}{2},\qquad
 \int_0^1\frac{t-t^2}{2}dt=\frac1{12}
\]
gives, in ordinary cohomological degrees $|m_n|=2-n$,
\[
 m_3(e_1J(z_1),\eta J(z_2),\eta J(z_3))
       =-\frac\eta{12}J(z_1)J(z_2)J(z_3).
\]
The degree-one output reflects the interval cochains.  It is compatible
with the strict degree-zero vertex algebra because this interval
resolution has nonzero differential.  It is not a cubic interaction
in the BF action.

The quantum boundary map above identifies actual supported
observables on the boundary with their bulk--boundary pushforward.
The current comparison identifies its polynomial chiral vacuum, and
the interval argument identifies the transferred separated products.
These three maps have different sources.  In particular, the last
map alone does not describe a bulk action, a collision-chain
compactification, or higher operations of an interacting theory.
The zero-level BF construction supplies a physical realization of
the zero-level current while retaining the additional bulk gauge and
field cohomology computed above.
